\section{Proof of Theorem~\ref{thm:hard}: the reduction from Grid Tiling}\label{app:hardness}
This appendix contains the construction and all proofs for Section~\ref{sec:hardness}. Steps that use a computer are marked \textbf{(CA1)}--\textbf{(CA4)}; they are listed in Section~\ref{sec:hardcomp} and collected in a status table at the end of the appendix. The appendix is organised so that each lemma can be corrected locally: Lemmas~\ref{lem:hchain}--\ref{lem:hprof} are statements about one chain; Lemma~\ref{lem:hblow} and Proposition~\ref{prop:hTR} give the ZE semantics; Lemmas~\ref{lem:hsplit}, \ref{lem:hregion} and Theorem~\ref{thm:hS} give soundness; Lemmas~\ref{lem:hpar}--\ref{lem:hpm} the parameters.

\subsection{Conventions}\label{sec:hconv}
Soft and hard pairs, members, states and zero-excess (ZE) solutions are as in Section~\ref{sec:prelim} and Appendix~\ref{app:tunnels}. A \emph{ZE partition} of an instance $(G,H)$ is a partition of $V(G)$ whose cost is $|H|$; by Lemma~\ref{lem:decomp} and Lemma~\ref{lem:tZE} it is a partition in which no hard pair is edited (hard edges lie inside clusters, hard non-edges join different clusters) and every member has exactly one edited pair. For a member $(a,b,c)$ with centre $b$ the state is $L$ if the single edit deletes $ab$ ($b\sim c$, $a$ apart), $R$ if it deletes $bc$, and $M$ if it inserts $ac$; these are the states $1,2,3$ of Appendix~\ref{app:tunnels}. Write $x\sim y$ for ``$x,y$ are in one cluster''. The \emph{excess} of a solution $S$ is $\mathrm{ex}(S)=|S|-|H|$. Grid Tiling instances are assumed to have all $S_{ij}\ne\emptyset$ and $n\ge2$ (an instance with an empty $S_{ij}$ is mapped to a fixed NO instance; $n=1$ is trivial).

\subsection{The thin chain}\label{sec:hthin}
\paragraph{Definition.} Fix $m\ge7$. The \emph{thin chain} $\mathrm{OH}_m$ has the vertices $s_0,\dots,s_{2m}$ and $e_0,\dots,e_m$ and the following members, listed as $(a,b,c)$ with centre $b$:
\begin{itemize}
\item spine members $P_j=(s_{2j},s_{2j+1},s_{2j+2})$, $0\le j\le m-1$;
\item $G_j=(e_j,s_{2j+1},e_{j+2})$ for $2\le j\le m-4$ and $F_j=(e_j,s_{2j+2},e_{j+1})$ for $2\le j\le m-3$;
\item left cap: $(e_0,s_2,e_1)$, $(s_1,e_1,s_3)$, $(e_0,s_4,e_2)$;
\item right cap: $(s_{2m-5},e_{m-3},s_{2m-1})$, $(e_{m-2},s_{2m-2},e_m)$, $(s_{2m-4},e_{m-1},e_m)$, $(s_{2m-3},e_{m-1},s_{2m})$, $(e_{m-2},e_{m-1},s_{2m-1})$, $(s_{2m-3},e_m,s_{2m-1})$.
\end{itemize}
The edges are the pairs $ab,bc$ of the members and the four \emph{hard edges} $s_0e_1$, $e_0s_3$, $s_{2m-3}e_{m-2}$, $s_{2m}e_m$; every other pair of vertices is a hard non-edge. The members have pairwise disjoint pair sets and are induced $P_3$'s (checked by computer, and by the pair-disjointness argument of Lemma~\ref{lem:hpm} for the thick version). The \emph{cell} of a vertex is $\lfloor t/2\rfloor$ for $s_t$ and $t$ for $e_t$; every member has its vertices in at most three consecutive cells, and every soft pair joins vertices in cells at distance at most $2$. A \emph{position} is an integer $a\in\{0,\dots,m-1\}$; a \emph{group} is a set $\{2u,2u+1\}$, $0\le u\le U-1$, where $m=2U$.

\paragraph{The partitions $P_a$.} Put $\varepsilon(0)=e_1$, $\varepsilon(1)=e_0$, $\varepsilon(j)=e_j$ for $2\le j\le m-2$ and $\varepsilon(m-1)=e_m$. For $0\le a\le m-1$ let $P_a$ be the partition of $V(\mathrm{OH}_m)$ whose clusters $C_0,\dots,C_{m-1}$ are
\[
C_j=\begin{cases}\{s_{2j},s_{2j+1}\}&j<a\\ \{s_{2j},s_{2j+1},s_{2j+2}\}&j=a\\ \{s_{2j+1},s_{2j+2}\}&j>a\end{cases}\ \cup\ \{\varepsilon(j)\}\ \ (\text{for }j\le m-2),
\]
together with the following changes at the right end. If $a\le m-2$ then
\[C_{m-1}=\{s_{2m-1},s_{2m},e_{m-1},e_m\};\]
if $a=m-1$ then $C_{m-2}$ contains in addition $e_{m-1}$ and $C_{m-1}=\{s_{2m-2},s_{2m-1},s_{2m},e_m\}$. In $P_a$ the spine members $P_0,\dots,P_{m-1}$ are in the states $R^{a}ML^{m-1-a}$.

\begin{lemma}[Chain lemma]\label{lem:hchain}
For every $m\ge7$ the ZE partitions of $\mathrm{OH}_m$ are exactly $P_0,\dots,P_{m-1}$; they are pairwise distinct.
\end{lemma}
\begin{proof}
\emph{(a) Each $P_a$ is a ZE partition.} This is a conjunction of local conditions: every member, and every hard pair that is an edge, involves vertices of at most three consecutive cells, and the pairs of vertices that must be apart are only constrained through clusters, which have vertices in at most three consecutive cells. A condition at a cell therefore depends only on the relation of the position $a$ to this cell (smaller by at least $3$, or $a-j\in\{2,1,0,-1,-2\}$, or larger), on the distance to the left end up to a constant, and on the distance to the right end up to a constant. Windows at the left end do not depend on $m$, windows at the right end depend only on the distance to the right end, and interior windows are isomorphic. All window types occur for $m\le 20$. \textbf{(CA1)}: for every $7\le m\le60$ and every $a$ the partition $P_a$ was checked by computer to have cost exactly $|H|$ with one edit per member and no hard pair edited (no failure). By the window argument this covers all $m\ge7$.

\emph{(b) There are no other ZE partitions.} \textbf{(CA2)}. Process the cells $0,1,\dots,m$ from left to right. After cell $c$, a \emph{partial ZE partition} is a partition of the vertices of the cells $\le c$ in which every hard pair and every member with all its vertices in these cells is respected. A vertex of cell $c'$ can be together with a vertex of a cell $c''\ge c'+3$ only if the pair is soft or a hard edge, which never happens (soft pairs join cells at distance at most $2$ and the hard edges join vertices in cells at distance at most $1$). So only the partition induced on the vertices of the last two cells matters for the future, together with the information which blocks already contain an older vertex (such a block is \emph{closed}). The state after cell $c$ is the set partition of the vertices of the cells $c-1,c$ with its set of closed blocks, with blocks renamed canonically and vertices described by (type, cell offset). At every interior cell there are exactly six reachable states, and the transition matrix between consecutive interior cells is the same matrix
\[
M=\begin{pmatrix}0&0&0&0&1&1\\0&0&0&0&0&0\\1&1&1&1&0&0\\0&0&0&0&0&0\\0&0&0&0&1&1\\0&0&0&0&0&0\end{pmatrix},\qquad\chi_M(x)=x^4(x-1)^2 ,
\]
with a Jordan block of size two at the eigenvalue $1$ (computed with exact integer arithmetic; the equality of the matrices was checked at the cells $7,\dots,19$ of $m=26$, and the transition rule at a cell involves only members and hard pairs within three consecutive cells, while the member list is translation invariant in the interior, so the matrix is the same at every interior cell of every $m$). The number $N_m$ of ZE partitions is $N_m=u^{\top}M^{\,m-c_0}v_m$ where $u$ is the vector after the left cap, which does not depend on $m$, and $v_m$ is determined by the right cap and depends only on the distance to the right end, hence is the same for all $m$ after the shift by the exponent. Since $x^4$ divides $\chi_M$, the generalised eigenspace of the eigenvalue $0$ is annihilated by $M^4$, and $M^{t}$ is for $t\ge4$ an affine function of $t$ (the Jordan block at $1$). Hence $N_m$ is an affine function of $m$ for $m\ge m_1$. A direct computation gives $N_m=m$ for $7\le m\le40$, so $N_m=m$ for all $m\ge7$. Together with (a), which provides $m$ distinct ZE partitions (their spine words $R^aML^{m-1-a}$ differ), these are all of them.
\end{proof}
\emph{Further evidence} (not used in the proof): an independent implementation (own validity checker, own SAT enumeration of all ZE partitions) found exactly the $m$ partitions $P_a$ for $m=7,\dots,14$ and for a sample of $m$ up to $100$; a third enumeration (SAT, all ZE partitions of the thin chain) gave exactly $m$ partitions for $m=14,18,26,42$ (the values of $m$ for $n=2,3,4,5$), with spine words $R^aML^{m-1-a}$.

\paragraph{Translation invariance, precisely.} The member list of $\mathrm{OH}_m$ consists of a \emph{left cap} (the members on the first few cells), a \emph{right cap} (the members on the last few cells) and the interior families $P_j,G_j,F_j$. The map $s_t\mapsto s_{t+2}$, $e_t\mapsto e_{t+1}$ shifts every cell by one and is an isomorphism between the sub-instances induced on the cells $\{j,j+1,j+2\}$ and $\{j+1,j+2,j+3\}$ whenever both windows avoid the caps. Consequently: (i) the constraints at the cells of the left cap involve only the cap members and the first interior members, the same for every $m$, so \emph{left cap windows are independent of $m$}; (ii) the constraints near the right end involve the right cap and the last interior members, the same for all $m$ after the shift of the index by $m$, so \emph{right cap windows depend only on the distance to the right end}; (iii) every other window is interior and isomorphic to every other interior window. In (CA1) every window type of (i)--(iii) occurs for $m\le20$, so a check of all $m\in[7,60]$ is a check of every window type, and the set of window types does not grow for larger $m$. In (CA2) the vector $u$ after the left cap is independent of $m$, the interior transfer matrix is the same $M$ at every interior cell (by the isomorphism above; checked at the cells $7,\dots,19$ of $m=26$), and the right-cap vector $v$ depends only on the distance to the right end, so $N_m=u^{\top}M^{\,m-c_0}v$ with $u,v,c_0$ fixed. The finite computation and the affine form of $M^t$ then determine $N_m$ for all $m\ge7$. This is a finite certificate with the invariance argument above, not a closed-form proof; the enumerations above are additional evidence.

\paragraph{Profiles.} Let $m=2U$. The \emph{host} of the group $u$ is $h_u=s_{4u+2}$; it is the last vertex of $P_{2u}$ and the first vertex of $P_{2u+1}$. Write $\varepsilon^-_u=\varepsilon(2u)$ and $\varepsilon^+_u=\varepsilon(2u+1)$; $\varepsilon^-_u\ne\varepsilon^+_u$ for all $u$.
\begin{lemma}[Profile of a host]\label{lem:hprof}
For the group $u$ the positions fall into four classes $A_0=\{a\le2u-1\}$, $A_1=\{2u\}$, $A_2=\{2u+1\}$, $A_3=\{a\ge 2u+2\}$, and for $a\in A_i$ the cluster of $h_u$ in $P_a$ is $\{h_u\}\cup\Pi_i$ with
\[
\Pi_0=\{s_{4u+1},\varepsilon^-_u\},\quad \Pi_1=\{s_{4u},s_{4u+1},\varepsilon^-_u\},\quad \Pi_2=\{s_{4u+3},s_{4u+4},\varepsilon^+_u\},\quad \Pi_3=\{s_{4u+3},\varepsilon^+_u\}.
\]
($A_0=\emptyset$ for $u=0$ and $A_3=\emptyset$ for $u=U-1$.) In particular $\Pi_0\subset\Pi_1$, $\Pi_3\subset\Pi_2$, $s_{4u}\in\Pi_1\setminus\Pi_0$, $s_{4u+4}\in\Pi_2\setminus\Pi_3$, and $\Pi_0,\Pi_1$ are disjoint from $\Pi_2,\Pi_3$.
\end{lemma}
\begin{proof}
Read off the definition of $P_a$: the vertex $s_{4u+2}$ is $s_{2j+2}$ for $j=2u$ and $s_{2j}$ for $j=2u+1$. If $a<2u$ the cluster $C_{2u}$ has the form $\{s_{4u+1},s_{4u+2}\}\cup\{\varepsilon(2u)\}$ and $C_{2u+1}\not\ni s_{4u+2}$; if $a=2u$ the cluster is $C_{2u}=\{s_{4u},s_{4u+1},s_{4u+2},\varepsilon(2u)\}$; if $a=2u+1$ it is $C_{2u+1}=\{s_{4u+2},s_{4u+3},s_{4u+4},\varepsilon(2u+1)\}$; if $a\ge2u+2$ it is $C_{2u+1}=\{s_{4u+2},s_{4u+3},\varepsilon(2u+1)\}$. At the caps, the label $\varepsilon$ takes the exceptional values stated before Lemma~\ref{lem:hchain}. \textbf{(CA3)}: the statement was in addition verified by computer for all even $m$ with $8\le m\le60$ and all $u$ ($1{,}836$ cases), including the groups $u=0$ and $u=U-1$.
\end{proof}
\paragraph{Allowed sets and forbidden positions.} For a set $W$ of thin vertices of one chain with $h_u\in W$ let $F_u(W)=\{a:\ \text{the cluster of }h_u\text{ in }P_a\not\subseteq W\}$ be the set of positions \emph{forbidden} by $W$. Define, for the group $u$,
\[
\mathrm{Pad}_u=\{h_u\}\cup\Pi_0\cup\Pi_3,\qquad X_u=\{h_u\}\cup\Pi_0,\qquad Y_u=\{h_u\}\cup\Pi_3,\qquad \mathrm{All}_u=\{h_u\}\cup\Pi_0\cup\Pi_1\cup\Pi_2\cup\Pi_3 .
\]
\begin{corollary}\label{cor:hforb}
$F_u(\mathrm{Pad}_u)=\{2u,2u+1\}$, $F_u(X_u)=\{a\ge2u\}$, $F_u(Y_u)=\{a\le2u+1\}$ and $F_u(\mathrm{All}_u)=\emptyset$.
\end{corollary}
\begin{proof}
By Lemma~\ref{lem:hprof}: for $\mathrm{Pad}_u$ the clusters at $A_0$ and $A_3$ are inside $\mathrm{Pad}_u$, while at $A_1$ the vertex $s_{4u}\notin\mathrm{Pad}_u$ and at $A_2$ the vertex $s_{4u+4}\notin\mathrm{Pad}_u$. For $X_u$: $A_0$ is allowed, $A_1$ is forbidden by $s_{4u}$, and $A_2,A_3$ by the non-empty sets $\Pi_2,\Pi_3$ disjoint from $\Pi_0$ (if $A_3=\emptyset$ there is nothing to forbid). $Y_u$ is symmetric: $A_3$ is allowed, $A_2$ is forbidden by $s_{4u+4}$, and $A_1,A_0$ by $\Pi_1,\Pi_0$, which are disjoint from $\Pi_3$. $\mathrm{All}_u$ contains every cluster.
\end{proof}

\subsection{Thick chains and guests}\label{sec:hthick}
Let $p\ge6$ be even. A fixed \emph{one-factorisation} $\mathcal F_1,\dots,\mathcal F_{p-1}$ of $K_p$ on the vertex set $\mathbb Z_p$ (round robin) is used throughout.

\paragraph{Thick chain.} The thick version of a thin chain with tag $c$ has the vertices $(v,i)$, $v$ a thin vertex, $i\in\mathbb Z_p$; $K_v=\{(v,i):i\in\mathbb Z_p\}$ is the \emph{super-vertex} of $v$. For each thin member $(a,b,c')$ there are the $p^2$ \emph{Latin members} $((a,i),(b,j),(c',i+j))$ with centre $(b,j)$, $i,j\in\mathbb Z_p$. For each thin hard pair $uv$ all $p^2$ pairs between $K_u$ and $K_v$ are hard, edges or non-edges as $uv$. Inside each $K_v$ all pairs are hard edges (the \emph{glue}) except the pairs used as pad or host pairs below. Every pair between $K_a$ and $K_b$ ($ab$ a pair of a thin member) is covered by exactly one Latin member: the pair $(a,i),(b,j)$ by the member with these two vertices, the pair $(b,j),(c',k)$ by $i=k-j$, and the pair $(a,i),(c',k)$ by $j=k-i$. So the Latin members have pairwise disjoint pair sets, and if $i+j\ne k$ the triple $(a,i),(b,j),(c',k)$ is a \emph{soft triangle}.

\paragraph{Guests.} A \emph{guest} is a single vertex $y$ together with a set $W_y$ of thin vertices (of one or two chains) such that all pairs between $y$ and the super-vertices $K_v$, $v\in W_y$, are soft and all other pairs of $y$ are hard non-edges, apart from the single-vertex members specified below. Realisation: for every $v\in W_y$ a perfect matching $\mathcal M(y,v)$ of $K_p$ is chosen, a different class $\mathcal F_r$ for different guests at the same $v$; for every $\{i,k\}\in\mathcal M(y,v)$ ($i<k$) there is the \emph{pad member} $(y,(v,i),(v,k))$ with centre $(v,i)$ (edges $y(v,i)$ and $(v,i)(v,k)$, non-edge $y(v,k)$). A \emph{permanent guest} $z$ with a \emph{host} $h\in W_z$ has instead, for the matching $\mathcal M(z,h)$, the \emph{host members} $((h,i),z,(h,k))$ with centre $z$ (edges $z(h,i)$, $z(h,k)$, non-edge $(h,i)(h,k)$), and pad members for all $v\in W_z\setminus\{h\}$. The pairs of $K_v$ in a matching used at $v$ are soft; the glue is the complete graph on $K_v$ minus these pairs. All pairs between different super-vertices that are not soft and not thin hard edges are hard non-edges, and so are all pairs between guests (except those in the members $(x,y,z)$ and $(z,y,z')$ below).

\paragraph{Thin-guest model.} A \emph{thin configuration} is a partition $\pi$ of the set of thin vertices, guests, anchors and stubs of the instance such that: (T1) a pair that is a hard edge in the thin picture (a thin hard edge, an uncut tunnel, a pair $(z,h)$ of a permanent guest and its host) lies inside a cluster; every other pair of two elements is split, except the \emph{free pairs} $(y,v)$, $v\in W_y$, and the pairs of the members; (T2) every thin chain member and every single-vertex member $(x,y,z)$, $(z,y,z')$ has exactly one edited pair.

\begin{lemma}[Blow-up]\label{lem:hblow}
Assume that the glue graph of every $K_v$ is connected (Lemma~\ref{lem:hpar}). Then the ZE partitions of a thick instance are exactly the lifts (replace each thin vertex $v$ by $K_v$) of the thin configurations. In particular in a ZE partition (a) every $K_v$ lies in one cluster; (b) a guest $y$ is together with $K_v$ only if $v\in W_y$; (c) a permanent guest is together with its host $K_h$; (d) the thick members of a thin member are all in the state of the thin member.
\end{lemma}
\begin{proof}
Let $\sigma$ be a ZE partition. No hard pair is edited, so the connected glue puts each $K_v$ into one cluster, which is (a). A pair between $K_u$ and $K_v$ that is a thin hard pair is hard of the same type for all $p^2$ pairs, so $K_u\sim K_v$ iff the thin pair is an edge. For a thin member $(a,b,c')$ the $p^2$ Latin members see the same togetherness pattern of the three super-vertices, so they are in the same state, which is a state of the thin member (one edit); this is (d). A pad member $(y,(v,i),(v,k))$ with $(v,i)\sim(v,k)$ has the togetherness patterns $(1,1,1)$ (if $y\sim K_v$) with one edit, inserting $y(v,k)$, and $(0,1,0)$ (if $y\not\sim K_v$) with one edit, deleting $y(v,i)$; so a pad member imposes no condition. If $v\notin W_y$ the pairs $y(v,i)$ are hard non-edges, so $y\not\sim K_v$, which is (b). A host member $((h,i),z,(h,k))$ with $(h,i)\sim(h,k)$ has one edit iff $z\sim(h,i)$ (pattern $(1,1,1)$: insert $(h,i)(h,k)$); if $z$ is apart it has two edits. This is (c). Conversely, the lift of a thin configuration has no edited hard pair (hard pairs between super-vertices and glue edges are respected by (T1), pairs of guests with $K_v$, $v\notin W_y$, are split) and one edit in every thick member, by the same case analysis. Note that in a lift the host pair $(h,i)(h,k)$ is a soft non-edge inside a cluster and is edited, while a pad pair is a soft edge inside a cluster and is not.
\end{proof}

\subsection{The reduction}\label{sec:hred}
Let $(k,n,(S_{ij}))$ be given, $n\ge2$, and let $p$ be the least even integer $\ge k^2+5$.

\paragraph{Chains and groups.} For every cell $(i,j)\in[k]^2$ there are two chains $A_{ij}$ and $B_{ij}$, each a thick copy of $\mathrm{OH}_m$ with $m=2U$, $U=1+n(n+1)$ (so $m\ge14$). For $\alpha\in[n]$ the \emph{block} $\alpha$ has the \emph{slot groups} $u=(\alpha-1)(n+1)+s$, $s\in[n]$, and is followed by the \emph{separator group} $u_\alpha=\alpha(n+1)$; group $0$ is the left end pad, and the last separator group $u_n=U-1$ is the right end pad.

\paragraph{Tiles and tunnels.} A tile $(a,b)\in S_{ij}$ has its \emph{slot} in $A_{ij}$ at the group $(a-1)(n+1)+b$ (block $a$) and in $B_{ij}$ at the group $(b-1)(n+1)+a$ (block $b$). The groups that are slots of tiles are \emph{tile groups}.

\paragraph{Guests.} For every chain $c$ and every group $u$ there is a permanent guest $z_{c,u}$ with host $h_u$; if $u$ is a tile group, $W_z=\mathrm{All}_u$, and in addition there are a \emph{lock guest} $y_{c,u}$ with $W_y=\mathrm{Pad}_u$, a vertex $x_{c,u}$ (the \emph{anchor}) and the single-vertex member $(x_{c,u},y_{c,u},z_{c,u})$ with centre $y_{c,u}$. If $u$ is not a tile group, $W_z=\mathrm{Pad}_u$ (a \emph{permanent lock}). All other pairs of $x,y,z$ are hard non-edges.

\paragraph{Tunnels.} For every cell and every tile $e\in S_{ij}$ the two anchors of $e$ (in $A_{ij}$ and in $B_{ij}$) are joined by a hard edge, the \emph{tunnel} of $e$. The anchors have no other hard edge.

\paragraph{Comparators.} For $\alpha\in[n-1]$, for every row $i$ and $j<k$, let $c_1=A_{ij}$, $c_2=A_{i,j+1}$, and for every column $j$ and $i<k$ let $c_1=B_{ij}$, $c_2=B_{i+1,j}$. In each case there is a \emph{comparator guest} $y$ with $W_y=X_{u_\alpha}(c_1)\cup Y_{u_\alpha}(c_2)$ (sets of thin vertices of the two chains) and the single-vertex member $(z_{c_1,u_\alpha},y,z_{c_2,u_\alpha})$ with centre $y$.

\paragraph{The instance.} $G$ is defined by the members above (edges $ab,bc$ of every member), the hard edges (thin hard edges, glue, tunnels) and no other edges. $H$ is the set of all members: Latin members, host members, pad members, lock members $(x,y,z)$ and comparator members. The excess is $\ell=k^2$. The instance is denoted $I_{\mathrm{tile}}$; for a set $C$ of tunnels, $I_C$ is $I_{\mathrm{tile}}$ with the pairs of $C$ made non-edges (they stay hard). A \emph{stub instance} is the same with some anchors' tunnels ending at stubs (vertices of degree one).

Lemma~\ref{lem:hpar} below shows that $H$ is a modification-disjoint packing of induced $P_3$'s, that the glue graphs are connected and that matchings suffice, so Lemma~\ref{lem:hblow} applies to $I_{\mathrm{tile}}$, to every $I_C$, and to the stub instances used below.

\subsection{Tunnel-release semantics}\label{sec:hTR}
For a chain $c$ and a cell let $\beta(c)\in[n]$ denote the \emph{block} of the tile whose tunnel is cut: $\beta(A_{ij})=a$ and $\beta(B_{ij})=b$ if the tile is $(a,b)$.

\begin{proposition}[Tunnel release]\label{prop:hTR}
Let $t_{ij}=(a_{ij},b_{ij})\in S_{ij}$ be a tile of every cell, $C$ the set of their tunnels, and $\sigma$ a ZE partition of $I_C$. Then in $\sigma$ every chain $c$ is in a position $a_c$ that lies in the group of its cut tunnel, and $a_{ij}\le a_{i,j+1}$, $b_{ij}\le b_{i+1,j}$ for all cells. Conversely, if $(t_{ij})$ satisfies these inequalities, $I_C$ has a ZE partition.
\end{proposition}
\begin{proof}
($\Rightarrow$) By Lemma~\ref{lem:hblow} take the thin configuration. Its restriction to a chain $c$ is a ZE partition of $\mathrm{OH}_m$ (all non-soft pairs inside the chain are hard), so by Lemma~\ref{lem:hchain} it is $P_{a_c}$ for one position $a_c$.

\emph{Permanent guests.} By Lemma~\ref{lem:hblow}(c), $z=z_{c,u}\sim K_{h_u}$, so $z$ is in the cluster of $h_u$ in $P_{a_c}$; for each $v$ in that cluster the pair $zK_v$ is not a hard non-edge, so $v\in W_z$ by (b). Hence $a_c\notin F_u(W_z)$. If $u$ is not a tile group, Corollary~\ref{cor:hforb} gives $a_c\notin\{2u,2u+1\}$.

\emph{Locks.} Let $u$ be a tile group whose tunnel $xx'$ ($x=x_{c,u}$) is not cut. It is a hard edge of $I_C$ with $N_h(x)=\{x'\}$ and $Q(x)\cap Q(x')=\emptyset$ (the soft partners are guests of different chains), so by Lemma~\ref{lem:tlock} the cluster of $x$ is $\{x,x'\}$; in particular $x\not\sim y,z$ for the members $(x,y,z)$. In the member $(x,y,z)$ with centre $y$ this leaves the single edit ``delete $xy$'' only if $y\sim z$ (if $y\not\sim z$ also $yz$ is deleted: two edits). So $y\sim z$ and $y$ is in the cluster of $h_u$; as above $a_c\notin F_u(\mathrm{Pad}_u)=\{2u,2u+1\}$.

Exactly one tunnel of the chain $c$ is cut (the tunnel of $t_{ij}$ serves both $A_{ij}$ and $B_{ij}$). All groups except the group of that tunnel exclude their two positions, and the groups partition $\{0,\dots,m-1\}$; hence $a_c$ lies in the group of the cut tunnel. With the slot indices above, this gives for the group $u_c$ of $c$: $u_c=(a_{ij}-1)(n+1)+b_{ij}$ for $A_{ij}$ and $(b_{ij}-1)(n+1)+a_{ij}$ for $B_{ij}$. Since $s\in[n]$, the group $u_c$ lies in the block $\beta(c)$: $u_c<u_\alpha$ iff $\beta(c)\le\alpha$, and $u_c>u_\alpha$ iff $\beta(c)\ge\alpha+1$; so $a_c\le2u_\alpha-1$ iff $\beta(c)\le\alpha$ and $a_c\ge2u_\alpha+2$ iff $\beta(c)\ge\alpha+1$.

\emph{Comparators.} Let $y$ be the comparator guest of $(c_1,c_2,\alpha)$, $u=u_\alpha$. In the member $(z_1,y,z_2)$ ($z_i=z_{c_i,u}$) the state $M$ would give $z_1\sim z_2$, impossible because $z_1,z_2$ lie in clusters of different chains (a pair of thick vertices of different chains is a hard non-edge). So the state is $L$ ($y\sim z_1$) or $R$ ($y\sim z_2$). In the first case the cluster of $h_u(c_1)$ is inside $W_y$, i.e.\ inside $X_u(c_1)$ (the part of $W_y$ in $c_1$), so $a_{c_1}\notin F_u(X_u)$, i.e.\ $a_{c_1}\le2u-1$, i.e.\ $\beta(c_1)\le\alpha$. In the second case $a_{c_2}\notin F_u(Y_u)$, i.e.\ $a_{c_2}\ge2u+2$, i.e.\ $\beta(c_2)\ge\alpha+1$. So for every $\alpha\in[n-1]$: $\beta(c_1)\le\alpha$ or $\beta(c_2)>\alpha$, which is equivalent to $\beta(c_1)\le\beta(c_2)$ (if $\beta(c_1)>\beta(c_2)$ take $\alpha=\beta(c_2)\le n-1$). For rows, $\beta(A_{ij})=a_{ij}$, for columns $\beta(B_{ij})=b_{ij}$, which are the inequalities of Grid Tiling.

($\Leftarrow$) Put every chain $c$ at the position $a_c=2u_c$, $u_c$ the group of its cut tunnel. Build the thin configuration: the base clusters are those of $P_{a_c}$. \emph{Permanent guests:} $z_{c,u}$ joins the cluster of $h_u$; this is allowed: for a tile group the set is $\mathrm{All}_u$, for the others $a_c\notin\{2u,2u+1\}$ and Corollary~\ref{cor:hforb} applies. \emph{Locks:} for a tile group $u\ne u_c$ (tunnel not cut) the lock guest $y_{c,u}$ joins the cluster of $z_{c,u}$ (allowed, $a_c\notin\{2u,2u+1\}$), and the anchors $x,x'$ at the two ends of that tunnel form the cluster $\{x,x'\}$; the member $(x,y,z)$ has the single edit ``delete $xy$'' (state $L$). For $u=u_c$ the tunnel is cut, $y_{c,u}$ is not joined to anything but its anchor, and the anchor forms the cluster $\{x,y\}$ (state $R$: delete $yz$). The pair $xx'$ is a cut tunnel, a hard non-edge, and the two anchors are in different clusters. \emph{Comparators:} the guest $y$ of $(c_1,c_2,\alpha)$ joins $z_1$ if $\beta(c_1)\le\alpha$ (allowed: then $a_{c_1}\le2u_\alpha-1$, so the cluster of $h_{u_\alpha}$ is inside $X_{u_\alpha}$) and joins $z_2$ otherwise (then $\beta(c_2)>\alpha$ and $a_{c_2}\ge2u_\alpha+2$, inside $Y_{u_\alpha}$). Every guest is in at most one cluster with a base vertex and two guests are never together, except in the pairs $y,z$ of members: guests join base clusters only through a permanent guest; two hosts $h_u,h_{u'}$ are never in one cluster (their positions in the spine differ by at least four); a tile group carries one lock guest and no comparator, and a separator group carries up to two comparator guests on the same $z$ (one as first chain with condition $\beta\le\alpha$, one as second chain with $\beta\ge\alpha+1$), which cannot both join. All thin chain members have one edit since $P_{a_c}$ is ZE, the members $(x,y,z)$ and $(z_1,y,z_2)$ have one edit as described, and every remaining pair is as required in (T1). By Lemma~\ref{lem:hblow} this lifts to a ZE partition of $I_C$.
\end{proof}

\subsection{Soundness and completeness}\label{sec:hsound}
\begin{lemma}[Parameters and validity]\label{lem:hpar}
(a) $H$ is a modification-disjoint packing of induced $P_3$'s of $G$. (b) Let $u(v)$ be the number of matchings used at $K_v$ (pad matchings and the host matching). Then $u(v)\le3$ for all $v$ and $u(v)\le 2$ if $k\le2$; at a vertex lying in four thin members $u(v)\le2$, and at a vertex in at most three thin members $u(v)\le3$. (c) The glue graph of every $K_v$ is connected, and $p-1\ge3$ matchings are available. (d) $c(H)\le 4p+2$ and $|V(G)|=O(k^2n^2p)$, $|H|=O(k^2n^2p^2)$.
\end{lemma}
\begin{proof}
(a) Every pair between guests and super-vertices, and every pair inside a super-vertex that is used as a pad or host pair, lies in exactly one member because different guests use different classes of the one-factorisation at the same $v$ and a pair of a matching lies in exactly one pad (or host) member; Latin members are pair-disjoint as shown above and have no pair inside a super-vertex; the single-vertex members $(x,y,z)$, $(z,y,z')$ use pairs between guests and anchors only. Every member $(a,b,c)$ has the edges $ab,bc$ and the non-edge $ac$ by definition, and these adjacencies are consistent because a soft pair belongs to one member.

(b) \textbf{(CA4)} A hand count from Lemma~\ref{lem:hprof}: the sets $W$ contain $h_u$, $\Pi_0,\Pi_3$ (pad), $\Pi_0$ (X), $\Pi_3$ (Y) or all $\Pi_i$ (All) of a single group $u$, so a super-vertex is used only by guests of the group $u$, except $s_{4u+4}=s_{4(u+1)}$, which occurs in $\Pi_2$ of the group $u$ and $\Pi_1$ of the group $u+1$. A group carries: a tile group the permanent guest $z$ and a lock guest $y$ (both $W\ni h_u$); a separator group (with $W_z=\mathrm{Pad}_u$) the permanent guest and at most two comparator guests (one with the chain as first chain, using $X_u$, one as second chain, using $Y_u$); any other group only $z$. At the host $h_u$: the permanent guest $z$ (host matching) plus either the lock guest $y$ of a tile group (a tile group carries no comparator), giving $u=2$, or, at a separator group, the up to two comparator guests (one with $h_u$ in $X_u$ of the first chain, one in $Y_u$ of the second), giving $u\le3$; a group is a tile group, a separator group or neither, never two of these. At $s_{4u+1},\varepsilon^-_u$: $z$ plus $y$ or the comparator guest of the first chain: at most $2$; at $s_{4u+3},\varepsilon^+_u$: $z$ plus $y$ or the comparator of the second chain: at most $2$; at $s_{4u},s_{4u+4}$: only in $\mathrm{All}$, i.e.\ only the permanent guests of tile groups, at most one from each of two groups: at most $2$. If $k=1$ there are no comparators, and if $k=2$ each chain has at most one neighbour, so the host has at most $2$. So a host carries $2$ (host matching and lock-guest pad) or $3$ (host matching and two comparator pads at a separator) matchings, and every non-host super-vertex at most $2$. The value $3$ occurs, at a separator host, only for $k\ge3$; the count of $u_{\max}$ by computer (thin-guest model, all tiles present, $1\le k\le5$, $2\le n\le4$) gave $2$ for $k\le2$ and $3$ for $k\ge3$, never more. The thin member counts: a vertex lying in four thin members is an $e$-vertex, and the hosts $s_{4u+2}$ lie in three; this, and the above counts, were confirmed by computer for $k\le4$ and $n\le3$ (maximum of $u(v)$ and of the number of thin members through every vertex, with the exceptional values at the caps).

(c) A vertex of $K_v$ has glue degree at least $p-1-u(v)$. For $p\ge8$ and $u(v)\le3$ this is at least $p-4\ge(p-1)/2$, hence the glue graph is connected (diameter at most $2$). For $p=6$ we have $k=1$, so $u(v)\le2$ and the degree is at least $3\ge(p-1)/2$. At most $3\le p-1$ matchings are used at any vertex.

(d) A thick vertex $(v,i)$ lies in $p$ Latin members of each thin member through $v$, and in $u(v)$ guest members. By (b) this is at most $4p+2$ at a vertex in four thin members and at most $3p+3\le4p+2$ otherwise. The sizes: $2k^2$ chains with $O(n^2)$ cells each and $O(p)$ thick vertices per thin vertex; each thin member gives $p^2$ Latin members; guests contribute $O(p)$ members each.
\end{proof}

\begin{lemma}[No super-vertex is split]\label{lem:hsplit}
Let $S$ be a solution of $I_{\mathrm{tile}}$ with $\mathrm{ex}(S)\le k^2$. Then every $K_v$ lies in one cluster of $G\sd S$. Consequently no hard glue pair and no pad pair (a soft edge inside $K_v$; in particular no pad pair of a comparator pad member) is edited. It is \emph{not} true that no pair inside a super-vertex is edited: the host pairs $(h,i)(h,k)$ are soft non-edges inside $K_h$ and are edited (Lemma~\ref{lem:hblow}).
\end{lemma}
\begin{proof}
Suppose $K_v$ meets at least two clusters; let $X\subset K_v$ be its smallest part, $x=|X|\in[1,p/2]$, $Y=K_v\setminus X$. Every pair between $X$ and $Y$ that is an edge of $G$ is deleted. A vertex of $X$ has at most $u(v)\le3$ soft pairs inside $K_v$, so at least $p-x-u(v)$ hard edges to $Y$; altogether at least $\varphi(x)=x(p-x-u(v))$ hard edits. Since $\varphi$ is concave, $\varphi(x)\ge\min\{\varphi(1),\varphi(p/2)\}$. We have $\varphi(1)=p-1-u(v)\ge p-4\ge k^2+1$. For $k\ge3$: $\varphi(p/2)\ge\frac p2(\frac p2-3)\ge\frac{(k^2+5)(k^2-1)}4>k^2$. For $k\le2$ we have $u(v)\le2$ and $p\in\{6,10\}$: $\varphi(p/2)\ge3>1$ resp.\ $\ge15>4$. So there are more than $k^2$ hard edits, but by Lemma~\ref{lem:decomp} the number of hard edits is at most $\mathrm{ex}(S)\le k^2$, a contradiction. Since $K_v$ lies in one cluster, every glue pair and every pad pair (all of them edges inside $K_v$) is an edge inside a cluster, hence not edited.
\end{proof}
(The host pairs $(h,i)(h,k)$ are soft \emph{non}-edges inside $K_h$; they are edited whenever the permanent guest $z$ joins its host, which happens in every ZE partition (state $M$, the single edit of the host member). What the proofs below use is only: each $K_v$ is inside one cluster, so no hard glue pair and no comparator pad pair is edited.)

\paragraph{Regions.} For a chain $c$ let $R_c$ be the set of its thick vertices and of its guests $z_{c,u},y_{c,u}$ and anchors $x_{c,u}$ (not the comparator guests). Let $T_c$ be the set of tunnel partners of anchors of $c$ (they are anchors of the partner chain, $A_{ij}\leftrightarrow B_{ij}$). The sets $R_c$ are pairwise disjoint, and $T_c\subseteq R_{c'}$ for the partner chain $c'$. Let $H_c$ be the set of members inside $R_c$ (Latin, host, pad members of its permanent and lock guests, lock members), $H_{\mathrm{cmp}}$ the set of members that contain a comparator guest (their pad members and the comparator members); $H$ is the disjoint union of the $H_c$ and $H_{\mathrm{cmp}}$. Let $\Pi_c$ be the set of pairs inside $R_c\cup T_c$ except the pairs with both ends in $T_c$. A \emph{cross-anchor pair} is a pair $\{x,x''\}$ of an anchor $x$ of $c$ carrying a tunnel and an anchor $x''$ of the partner chain carrying a tunnel; cross-anchor pairs are hard (tunnels and hard non-edges). A pair lies in two sets $\Pi_c,\Pi_{d}$ iff it is a cross-anchor pair (and then $d$ is the partner of $c$); every other pair lies in at most one $\Pi$. Pairs of $H_{\mathrm{cmp}}$ members other than the pad pairs involve a comparator guest or join $z_{c_1}$ to $z_{c_2}$ and lie in no $\Pi_c$.

\begin{lemma}[Region bound]\label{lem:hregion}
Let $S$ be a solution with $\mathrm{ex}(S)\le k^2$. For a chain $c$ and a member $M$ let $x_M=S\cap\mathrm{pairs}(M)$; let $h_c$ be the number of edits in $\Pi_c$ on pairs that belong to no member of $H_c$, let $E_c=\sum_{M\in H_c}(|x_M|-1)\ge0$ be the \emph{member excess} (members of $H_c$ with two or three edits), and put $r_c=|S\cap\Pi_c|-|H_c|$. Then $r_c=h_c+E_c\ge\max\{1,h_c\}$.
\end{lemma}
\begin{proof}
By Lemma~\ref{lem:hsplit} the pad pairs of comparator pad members (which lie in $\Pi_c$ but belong to members of $H_{\mathrm{cmp}}$, not of $H_c$) are unedited, so the edits in $\Pi_c$ are on pairs of members of $H_c$ or on pairs of no member at all; this gives $|S\cap\Pi_c|=\sum_{M\in H_c}|x_M|+h_c$, i.e.\ $r_c=h_c+E_c$. (The definition of $h_c$ must say ``no member of $H_c$'': a pad pair of a member of $H_{\mathrm{cmp}}$ lies in $\Pi_c$ but, being unedited, never counts.) Each member has at least one edit, so $E_c\ge0$ and $r_c\ge h_c$; if $h_c\ge1$ we are done. Suppose $h_c=0$ and $E_c=0$, i.e.\ $r_c=0$: then every member of $H_c$ has exactly one edit and $S$ has no other edit in $\Pi_c$. In particular no tunnel at $c$ is cut. Two stubs $x_1',x_2'\in T_c$ (partners of anchors $x_1,x_2$ of $c$) are not together: otherwise $x_1\sim x_1'\sim x_2'\sim x_2$ (tunnels unedited), but $x_1x_2$ is an unedited hard non-edge. So the partition restricted to $R_c\cup T_c$ has no edit on pairs with both ends in $T_c$ either, and is a ZE partition of the instance $G_c=G[R_c\cup T_c]$ with the packing $H_c$ (the comparator pad pairs inside $R_c$ are plain hard edges of $G_c$, and the vertices of $T_c$ are stubs: each has as only edge its tunnel). $G_c$ is a stub instance, so Lemma~\ref{lem:hblow} applies. The proof of Proposition~\ref{prop:hTR} ($\Rightarrow$) for the single chain $c$, with all tunnels present, uses only Lemma~\ref{lem:hchain}, Corollary~\ref{cor:hforb}, the permanent guests and the locks, and no comparators; it shows that the position $a_c$ lies outside every group (every group is excluded either by a permanent lock or by an uncut tunnel), a contradiction since the groups cover all positions. Hence $r_c\ge1$ whenever $h_c=0$, and $r_c\ge\max\{1,h_c\}$ in all cases.
\end{proof}

\begin{theorem}[Soundness]\label{thm:hS}
Every solution $S$ of $I_{\mathrm{tile}}$ has $\mathrm{ex}(S)\ge k^2$. If $\mathrm{ex}(S)=k^2$, then $S$ consists of one cut tunnel in every cell and one edit in every member, and nothing else; hence $S\setminus C$ is a ZE partition of $I_C$ with $C$ the set of cut tunnels, and Grid Tiling has a solution.
\end{theorem}
\begin{proof}
Let $\mathrm{ex}(S)\le k^2$ (otherwise there is nothing to show). Let $q$ be the number of edited cross-anchor pairs and $S_{\mathrm{rest}}$ the set of edits outside $\bigcup_c\Pi_c$. Since only cross-anchor pairs lie in two sets, $|S|=\sum_c|S\cap\Pi_c|-q+|S_{\mathrm{rest}}|$. Every member of $H_{\mathrm{cmp}}$ has an edit; its pad pair is unedited because no comparator pad pair is edited (Lemma~\ref{lem:hsplit}), and its other pairs are in no $\Pi_c$; these members are pair-disjoint, hence $|S_{\mathrm{rest}}|\ge|H_{\mathrm{cmp}}|$. With $r_c=|S\cap\Pi_c|-|H_c|$ (Lemma~\ref{lem:hregion}) and $|H|=\sum_c|H_c|+|H_{\mathrm{cmp}}|$,
\[
\mathrm{ex}(S)\ =\ \sum_c r_c-q+\bigl(|S_{\mathrm{rest}}|-|H_{\mathrm{cmp}}|\bigr),\qquad r_c=h_c+E_c\ge\max\{1,h_c\},\qquad |S_{\mathrm{rest}}|-|H_{\mathrm{cmp}}|\ge0 .
\]
All three slacks $E_c$, $r_c-\max\{1,h_c\}$ and $|S_{\mathrm{rest}}|-|H_{\mathrm{cmp}}|$ are non-negative, so $\mathrm{ex}(S)\ge\sum_c\max\{1,h_c\}-q$.
Let $Z=\{c:h_c=0\}$ and let $h_{\mathrm{in}}$ be the number of hard edits lying in exactly one $\Pi_c$. A cross-anchor edit counts in two chains, so $\sum_ch_c=h_{\mathrm{in}}+2q$, and $\sum_c\max\{1,h_c\}=|Z|+\sum_{c\notin Z}h_c=|Z|+h_{\mathrm{in}}+2q$. Hence $\mathrm{ex}(S)\ge|Z|+h_{\mathrm{in}}+q+\sum_cE_c$. The number of chains not in $Z$ is at most $\sum_ch_c$, so $|Z|\ge2k^2-h_{\mathrm{in}}-2q$, and therefore $\mathrm{ex}(S)\ge2k^2-q$. On the other hand every edited cross-anchor pair is a hard edit, so $q\le|U|\le\mathrm{ex}(S)$ (Lemma~\ref{lem:decomp}). Thus $\mathrm{ex}(S)\ge\max\{q,2k^2-q\}\ge k^2$.

Let $\mathrm{ex}(S)=k^2$. Then $q=k^2$ (as $\mathrm{ex}\ge q$ and $\mathrm{ex}\ge2k^2-q$ give $q=k^2$), and every inequality above is tight: from $\mathrm{ex}(S)\ge|Z|+h_{\mathrm{in}}+q+\sum_cE_c$ we get $Z=\emptyset$, $h_{\mathrm{in}}=0$ and $\sum_cE_c=0$, and from the first display $|S_{\mathrm{rest}}|=|H_{\mathrm{cmp}}|$. Every chain has $h_c\ge1$ and $\sum_ch_c=2q=2k^2$, so $h_c=1$ for all $2k^2$ chains: the only non-member edit at a chain is one cross-anchor pair. All hard edits are these $k^2$ pairs (as $|U|\le\mathrm{ex}=k^2=q$). Since $E_c=0$, all members of every $H_c$ have exactly one edit, and since $|S_{\mathrm{rest}}|=|H_{\mathrm{cmp}}|$ with one edit at least per member, $S_{\mathrm{rest}}$ has exactly one edit per member of $H_{\mathrm{cmp}}$. In a cell $(i,j)$ exactly one cross-anchor pair is edited (each chain has $h_c=1$ and sees only cross pairs of its own cell). It is a tunnel: suppose the edited pair is a non-edge $x_Ax_B$ of anchors of $A_{ij}$ and $B_{ij}$ whose tunnels are $x_Ax_B^{*}$ and $x_A^{*}x_B$ (with $x_B^*\ne x_B$). It is inserted, so $x_A\sim x_B$; the tunnel $x_Ax_B^*$ is an unedited edge, so $x_A\sim x_B^*$; hence two distinct anchors of $B_{ij}$ are together, but their pair is an unedited hard non-edge ($h_{B_{ij}}=1$ is used up by the cross pair), a contradiction. So the edited cross pair is a tunnel of the cell, deleted, and $S$ is the set $C$ of $k^2$ cut tunnels (one per cell) together with one edit per member. By Lemma~\ref{lem:hsplit} the partition is a partition with $K_v$ intact, and it is a ZE partition of $I_C$. Proposition~\ref{prop:hTR} ($\Rightarrow$) gives a tiling: the tiles $t_{ij}$ of the cut tunnels satisfy $a_{ij}\le a_{i,j+1}$ and $b_{ij}\le b_{i+1,j}$.
\end{proof}

\begin{theorem}[Completeness]\label{thm:hC}
If the Grid Tiling instance has a solution $(s_{ij})$, then $\OPT(I_{\mathrm{tile}})=|H|+k^2$.
\end{theorem}
\begin{proof}
Let $C$ be the tunnels of the tiles $s_{ij}$. By Proposition~\ref{prop:hTR} ($\Leftarrow$) the instance $I_C$ has a ZE partition, whose edit set (relative to $G$) is the set of its $|H|$ member edits together with the $k^2$ deleted tunnels. So $\OPT\le|H|+k^2$, and $\OPT\ge|H|+k^2$ by Theorem~\ref{thm:hS}.
\end{proof}

\subsection{Pathwidth, pair-maximality and proof of Theorem~\ref{thm:hard}}\label{sec:hpwpm}
\begin{lemma}[Pathwidth]\label{lem:hpw}
Every component of $I(H)$ is contained in the set of members of the $k$ chains of one row $\{A_{i,1},\dots,A_{i,k}\}$ or one column $\{B_{1,j},\dots,B_{k,j}\}$ together with their guests and comparator guests; so $I(H)$ has at most $2k$ components, and $\mathrm{pw}(I(H))<24kp^2$, independently of $n$.
\end{lemma}
\begin{proof}
Members of different chains share a vertex only through a comparator guest or through its pad members (they share vertices with Latin members of both chains); tunnels are hard pairs and are not members. This gives the components. Define a \emph{level} $\lambda(M)$ of a member: for a Latin member of a thin member, the smallest cell of the thin member (all three cells lie in $[\lambda,\lambda+2]$); for every member belonging to a guest of the group $u$ (host, pad, lock, comparator members), $\lambda=2u$. A guest of group $u$ has pad and host vertices in cells $[2u,2u+2]$ (the vertices $s_{4u},\dots,s_{4u+4},\varepsilon^\pm_u$), so for every thick vertex $(v,i)$ all members through it have levels in $[\mathrm{cell}(v)-2,\mathrm{cell}(v)]$, and all members through a guest have the same level. Hence two members that share a vertex have levels differing by at most $2$. For a component $Q$ put $B_j=\{M\in Q:\ j\le\lambda(M)\le j+2\}$; every member lies in the bags $j\in[\lambda(M)-2,\lambda(M)]$, an interval, and two members sharing a vertex have a common bag; so $(B_j)_j$ is a path decomposition. At one level a chain has at most $6$ thin members (at most $6$ at the right cap, at most $3$ elsewhere), hence at most $6p^2$ Latin members. At an even level $2u$, the guest members of one chain are at most $3.5p$ (permanent guest) $+\,2.5p+1$ (lock) in a tile group and $2.5p$ in a separator group, and the comparator guests of a row or column at a separator group contribute at most $k-1$ guests with at most $3p+1$ members each. So a level of a component has at most $6kp^2+k(6p+1)$ members, and $|B_j|\le3\cdot6kp^2+2k(6p+1)\le24kp^2$ for $p\ge6$. The bound does not depend on $n$ because the number of members at one level is a function of $p$ and $k$ only: the cell structure is the same in every chain, guests and comparators of a group sit at the same level $2u$ in every chain, and the level range of a member is at most $2$.
\end{proof}
\emph{Check of the decomposition} (not part of the proof). Independently, the path decomposition with bags $\{M:\lambda'(M)\le t\le\lambda''(M)\}$ (interval of the coordinates of the vertices of $M$) was verified by brute force: for every vertex of the thick graph the intervals of the members through it have a common point (Helly property of intervals), $0$ violations; every member spans at most $2$ cells. For $k=1$, $p=8$ and $n=2,3,4,5$ ($|V|=722,1314,2118,3118$) the largest bag has $1296,1296,1304,1296$ members, independent of $n$, as the proof states.

\begin{lemma}[Pair-maximality]\label{lem:hpm}
$H$ is pair-maximal.
\end{lemma}
\begin{proof}
An induced $P_3$ with three hard pairs consists of two hard edges sharing a vertex and a hard non-edge. The hard edges form the graph $\Gamma_h$ whose components are: (a) for a thin vertex without a thin hard edge, the glue of $K_v$; (b) for each of the four thin hard edges $uv$ of a chain, $K_u\cup K_v$ with the glue inside both and all edges between them; (c) the tunnels $xx'$ (the anchors carry no other hard edge, guests carry none). In the components (a) and (b) the pairs that are not adjacent are exactly the pad and host pairs inside a super-vertex, which are soft; a component (c) is a single edge. So in every component every non-adjacent pair is soft, and an induced $P_3$ with two hard edges has its three vertices in one component of $\Gamma_h$, hence its third pair is soft or an edge. Hence no induced $P_3$ has all pairs hard. (Checked by computer for $(k,n,p)=(1,2,8)$, $(1,2,6),\dots,(1,5,6)$ and $(2,2,10)$.)
\end{proof}

\begin{proof}[Proof of Theorem~\ref{thm:hard}]
The instance $(G,H,k^2)$ is constructed in time polynomial in $k$ and $n$ (Lemma~\ref{lem:hpar}(d)). (i) holds by definition, Lemma~\ref{lem:hpar}(a),(d) and Lemma~\ref{lem:hpm}; (ii) is Lemmas~\ref{lem:hpar}(d) and~\ref{lem:hpw} ($\tw\le\mathrm{pw}$); (iii) is Theorems~\ref{thm:hS} and~\ref{thm:hC} (a solution of size at most $|H|+k^2$ has excess at most $k^2$, hence exactly $k^2$, hence a tiling exists by Theorem~\ref{thm:hS}). The parameters $\ell+\tw(I(H))$, $\ell+\mathrm{pw}(I(H))$ and $\ell+c$ are bounded by $k^2+24k(k^2+6)^2$, which is a function of $k$ only, and the reduction is polynomial: it is a parameterized reduction from Grid Tiling with $\le$ parameterized by $k$, which is W[1]-hard. An algorithm of running time $f(\ell,w)n^{O(1)}$ would decide Grid Tiling in time $f'(k)\,n^{O(1)}$, so $\mathrm{FPT}=\mathrm{W[1]}$. The ETH statement is Corollary~\ref{cor:hardeth}.
\end{proof}

\begin{remark}[The chain object]\label{rem:hchaincore}
Consider one chain with $N$ tile groups, each with a lock guest and an anchor whose tunnel ends at a stub, all other groups being permanent locks. The proof of Proposition~\ref{prop:hTR} ($\Rightarrow$) shows that with all tunnels present there is no ZE partition, and the proof of ($\Leftarrow$) shows that cutting any single tunnel gives one. In the language of Appendix~\ref{app:tunnels} the set of all $N$ tunnels is therefore a core, of unbounded size $N$, while $c\le4p+2$ and the pathwidth do not depend on $N$ (the bound of Lemma~\ref{lem:hpw}). This is the family of tight instances with bounded $c$ and width and cores of unbounded size that the discussion of Appendix~\ref{app:tunnels} asks for; by Corollary~\ref{cor:sfcore} it must have soft triangles, and it does (the Latin members). The argument is the proof above applied to a single chain and was not run separately by computer. The hypotheses $p$ even, $p\ge8$ suffice for this statement (the tunnel semantics do not use the no-split lemma); for the full problem $p>\ell$ is needed.
\end{remark}

\subsection{Exact optimum on full thick instances}\label{sec:hexact}
This is evidence, not part of the proof. By Lemma~\ref{lem:hsplit}, every solution with $\mathrm{ex}\le k^2$ leaves each $K_v$ inside one cluster, so $\OPT\le|H|+k^2$ iff the minimum over solutions with every $K_v$ unsplit is at most $|H|+k^2$. That minimum is the optimum of a weighted correlation clustering instance on the \emph{contraction} (classes: the super-vertices and the single guest vertices; for two classes $A,B$ the weights are the numbers of edges and of non-edges of $G$ between them, computed from the real thick graph, so that every pad pair, host pair, thin member pair, hard non-edge and tunnel is counted; plus the internal non-edges of each class). All pairs of classes are editable and nothing is assumed about the states of the members; the problem is solved exactly by CP-SAT with lazily added transitivity constraints. The contraction was validated against the unrestricted exact solver (all pairs of $V(G)$ editable, no contraction) for $k=1$: with one tile $p=8$ both give $\OPT=|H|+1$ ($|H|=3258$), and without a tile ($p=4$) both give $\OPT=|H|+8$ ($\ell=0,\dots,7$ unsatisfiable, $\ell=8$ satisfiable, no split class). Results, with $\mathrm{ex}=\OPT-|H|$ (expected $\mathrm{ex}=k^2$ if Grid Tiling is solvable and $\mathrm{ex}\ge k^2+1$ otherwise):
\begin{center}
\small
\begin{tabular}{@{}cccllll@{}}
\hline
$k$ & $n$ & $p$ & instances & YES: $\mathrm{ex}$ & NO: $\mathrm{ex}$ & mismatches\\\hline
1 & 3 & 8 & 3 (random) & 3 YES, $\mathrm{ex}=1$ & -- & 0\\
2 & 2 & 10 & 11 (3 fixed, 8 random) & 6 YES, $\mathrm{ex}=4$ & 5 NO, $\mathrm{ex}=5$ & 0\\
2 & 3 & 10 & 2 (random) & 2 YES, $\mathrm{ex}=4$ & -- & 0\\
3 & 2 & 14 & 2 (fixed) & 1 YES, $\mathrm{ex}=9$ & 1 NO, $\mathrm{ex}=10$ & 0\\\hline
\end{tabular}
\end{center}
The largest instance has $|V(G)|=11{,}262$ and $|H|=153{,}498$ ($k=3$). On every instance solved, YES instances have $\mathrm{ex}=k^2$ exactly and NO instances have $\mathrm{ex}=k^2+1$ exactly (the cheapest deviation found by the solver is one comparator guest left alone, one member with two edits). That NO instances sit at exactly $k^2+1$ is an observation on these tests, not a theorem; the proof only uses $\mathrm{ex}\ge k^2+1$. This is the end-to-end exact test of $k\ge2$ that the unrestricted solver could not do; it uses Lemma~\ref{lem:hsplit}, which is proved by hand and was validated against the unrestricted solver on $k=1$.

\subsection{Status of the verification}\label{sec:hstatus}
\begin{center}
\small
\begin{tabular}{@{}p{0.34\textwidth}p{0.6\textwidth}@{}}
\hline
Statement & Status\\\hline
Lemma~\ref{lem:hchain}(a) ($P_a$ is ZE, all $m$) & proved modulo (CA1): finite check $m\le60$ plus window argument\\
Lemma~\ref{lem:hchain}(b) (no other ZE partition) & proved modulo (CA2): transfer matrix $\chi=x^4(x-1)^2$, $N_m=m$ for $7\le m\le40$, translation invariance spelled out above; independent SAT enumerations ($m=14,18,26,42$; sample up to $m=100$)\\
Lemma~\ref{lem:hprof}, Cor.~\ref{cor:hforb} (profiles) & read off the formula; (CA3) checked for even $m\le60$ and all groups\\
Lemma~\ref{lem:hblow} (blow-up) & proved by hand\\
Prop.~\ref{prop:hTR} (tunnel release, comparators) & proved by hand from the above; thin model tested on random instances ($k\le3$), thick comparator on all position pairs\\
Lemma~\ref{lem:hpar} ($u_{\max}\le3$, $c\le4p+2$) & hand count written out (host $2$ or $3$, non-host $\le2$); (CA4) confirmed by computer, $u_{\max}=2$ for $k\le2$, $3$ for $k\ge3$, $k\le5$\\
Lemma~\ref{lem:hsplit} (no split) & proved by hand (needs $p\ge k^2+5$); statement corrected: no hard glue pair and no pad pair is edited (host pairs are edited); validated on $k=1$ against the unrestricted exact solver\\
Lemma~\ref{lem:hregion}, Thm.~\ref{thm:hS} (soundness) & proved by hand, with the member-excess slack $E_c$ carried in the count (second audit); exact optimum on full thick instances (contraction of every $K_v$, justified by Lemma~\ref{lem:hsplit}; CP-SAT, all pairs editable), see Section~\ref{sec:hexact}: $k=1$, $n=3$, $p=8$: 3 YES ($\mathrm{ex}=1$); $k=2$, $n=2$, $p=10$: 6 YES ($\mathrm{ex}=4$), 5 NO ($\mathrm{ex}=5$); $k=2$, $n=3$, $p=10$: 2 YES ($\mathrm{ex}=4$); $k=3$, $n=2$, $p=14$ ($|V|=11{,}262$): 1 YES ($\mathrm{ex}=9$), 1 NO ($\mathrm{ex}=10$); contraction validated against the unrestricted solver for $k=1$; NO instances have $\mathrm{ex}=k^2+1$ in these tests (observation, not a theorem)\\
Thm.~\ref{thm:hC} (completeness) & proved by hand\\
Lemma~\ref{lem:hpw} (pathwidth $<24kp^2$) & proved by hand (level sweep); interval decomposition checked by brute force (Helly), largest bag independent of $n$\\
Lemma~\ref{lem:hpm} (pair-maximal) & proved by hand; checked on small instances\\\hline
\end{tabular}
\end{center}
